% =============================================================================
% Appendix D  The Published Circuits as Built  --  cited from
% sec:standard-published and sec:results-faithful.
% Plan: ThesisPlanning.md §6 (D).
%
% CONTENT  The paper's Figs 2, 8 and 9 rebuilt; every deviation from the printed
%          figure; the measured qubits, gates and eps_G(n); the published
%          prefactor against ||G||_2.
% SOURCES  Planning.md §3.7, §6.3, App. A-4; claims A-06 to A-08 (pihm
%          docs/18_claims.md: tests/claims/test_pillar1.py);
%          pihm/src/pihm/circuits/standard.py (published_derivative_encoding,
%          derivative_ladder_error_bound).
% KEPT from the draft: the ladder mechanism and its error bound (prop:epsilon-G)
%   -- pihm's derivative_ladder_error_bound states the same bound for the
%   Fig. 9 build. DROPPED (snapshot 2cd640c): the legacy reconstruction's
%   resource count (n + 4 ancilla, adder, "two flags live") and its
%   verification claims. The ladder's per-bit coefficient is written mu here
%   (the paper's alpha; the draft's beta): beta is imaginary time (T-08) and
%   alpha a subnormalisation.
% =============================================================================

\Cref{sec:pihm-encoding} describes, at the level a reader needs to follow the
argument, the circuit Wu et al.\ use to load $\G$'s column-linear entries as a
rotation amplitude \cite{wu2025pihm}. This appendix gives the mechanism in full,
together with a rigorous bound on its entrywise error that this thesis derives and
verifies -- the source paper states only that the error is ``negligible'' and
decreases ``exponentially'' with $n$, without a closed form.
% REVISE: "verifies" -- point to the approximation-mode measurement of
%   eps_G(n) (tab:verification-modes) at results-v1.

\subsection{The circuits as built}
\label{app:published-build}
% LAND: Figs 2 and 8 (feature map + zero-state reflection; B on n + 2 qubits,
%   alpha_B = sqrt(2^{n+1}); D0 with alpha divided by f(x_s)); Fig. 9 built as
%   drawn from the legible figure and the caption angles (the omega_c flag, the
%   X ancilla, the offset register with its Hadamard sandwich, the value
%   ancilla; U_R and U_C). EVERY deviation from the printed figure, one line
%   each. The subtractor/validity structure is shared with the structured
%   encoding (app:structured); only the value load differs.

\paragraph{The constraint rows (Figs 2 and 8).} $\Brow(x)$ and $\Dzero(x_s)$ are block-encoded by the inverse feature map $U_\tau^\dagger(x)$ followed by $\hat S_{n+1} = 2\ket{0}\!\bra{0}-I$ controlled inside a Hadamard sandwich on one ancilla, a two-term LCU on $n+2$ qubits. The encoded block is $\Brow(x)/\sqrt{2^{n+1}}$, and $\Dzero(x_s)/\bigl(\sqrt{2^{n+1}}/f(x_s)\bigr)$ for the collapse row. The global phase of $-1$ in $\hat S_{n+1}$ is a relative phase once the circuit is controlled, and cannot be dropped.

\paragraph{The derivative circuit (Fig.\ 9).} The registers are, in order, a flag qubit acted on by $R_Y(\omega_c)$ with $\omega_c = 2\arccos\sqrt{1/3}$, an $X$ ancilla, an offset register of $m = n$ qubits sandwiched between Hadamards, a value ancilla, and the $n$-qubit system. Two conditioned subroutines assign the entries on different bases. In $\hat U_R$ the offset $\ket{i}$ selects the $(2i+1)$th superdiagonal and the system $\ket{j}$ the $j$th column; it consists of $n$ singly-controlled $R_Y(\theta_k)$ gates on the value ancilla, controlled by the bits of $\ket{j}$, followed by $n$ cyclic shifts $\hat R_n$. In $\hat U_C$ the offset selects the $(i+1)$th row; it consists of $n$ multi-controlled $R_Y(\varphi_k)$ gates, one shift $\hat L_n$ and $2n$ swaps. With $\mu = 2/\|\G\|_S$ the angles are $\theta_k = 2\arcsin(-2^{k-1}\mu)$ and $\varphi_k = 2\arcsin[-2^{k-1}\mu(1/\sqrt2-1)]$. The circuit uses $m+n+3 = 2n+3$ qubits, its gate count is linear in $n$ before the multi-controlled gates are decomposed and $O(n^3)$ after an $O(n^2)$ decomposition of each, and $\G$ is the top-left $2^n\times2^n$ block of the unitary multiplied by $(2^{n-1}+2^n)\|\G\|_S$. These are the paper's statements; their tests follow.

\subsection{The small-angle ladder and its error}
\label{app:ladder}

\paragraph{The mechanism.} A ladder of $n$ singly-controlled $R_y$ rotations acts on
one flag qubit, initialised to $\ket{1}$: bit $m$ of the (register-arithmetic-shifted)
column index controls $R_y(\theta_m)$,
\begin{equation}
    \theta_m = 2\arcsin\!\big({-}2^m\mu\big), \qquad m = 0, \dots, n-1,
    \label{eq:app-ladder-primary}
\end{equation}
for a small real coefficient $\mu$. Because rotations on the same qubit compose by
exact angle addition ($R_y(a)R_y(b) = R_y(a+b)$, an $SO(2)$ identity), the net
rotation for a column index $p$ with bit set $S$ is $R_y(\Theta(p))$,
$\Theta(p) = \sum_{m \in S} \theta_m$, giving the flag qubit's $\ket{0}$-amplitude
\begin{equation}
    \mathrm{raw}(p) = -\sin\!\big(\Theta(p)/2\big) = \sin\!\Big(\textstyle\sum_{m \in S} \arcsin(2^m\mu)\Big).
    \label{eq:app-ladder-raw}
\end{equation}
The row-$0$ correction (entries $\sqrt{2}p$ in place of $2p$) is a second,
independently accumulated ladder on the same flag qubit, with per-bit coefficient
$\mu(1/\sqrt{2} - 1)$, controlled on the same ``this is row $0$'' flag the main
construction already computes.

\paragraph{Why this is only approximate.} $\mathrm{raw}(p)$ in \Cref{eq:app-ladder-raw}
would equal the desired linear value $\mu p$ exactly only if
$\sin\big(\sum_m \arcsin(x_m)\big) = \sum_m x_m$, which does not hold in general --
$\sin$ of a sum is not the sum of sines. The discrepancy is a genuine, quantifiable
small-angle error.

\begin{proposition}[Entrywise error of the ladder construction]
\label{prop:epsilon-G}
Let $J_3(p) := \sum_{m \in S} (2^m)^3$ for column index $p$ with bit set $S$. Then
\begin{equation}
    \big|\mathrm{raw}(p) - \mu p\big| \;\le\; \frac{\mu^3}{6}\,p^3 \;+\; O(\mu^5 N^5).
    \label{eq:app-ladder-bound}
\end{equation}
\end{proposition}

\begin{proof}
Expand both non-linearities to leading order in the small parameter $\mu$:
$\arcsin(y) = y + y^3/6 + O(y^5)$ and $\sin(x) = x - x^3/6 + O(x^5)$. Summing the
first expansion over the fired bits and substituting into the second,
\begin{align}
    \sum_{m \in S} \arcsin(2^m\mu) &= \mu p + \tfrac{\mu^3}{6} J_3(p) + O(\mu^5 N^5), \\
    \sin\!\Big(\sum_{m \in S} \arcsin(2^m\mu)\Big)
        &= \sin(\mu p) + \cos(\mu p)\cdot\tfrac{\mu^3}{6}J_3(p) + O(\mu^5N^5) \notag\\
        &= \mu p + \tfrac{\mu^3}{6}\big(J_3(p) - p^3\big) + O(\mu^5N^5).
\end{align}
Since $p = \sum_{m\in S} 2^m$ and every $2^m \le p$, the multinomial expansion of
$p^3$ term-by-term dominates $J_3(p)$, with equality only when $S$ is a single bit
($p$ itself a power of two): $J_3(p) \le p^3$. Taking the worst case $p = N-1$ gives
\Cref{eq:app-ladder-bound}.
\end{proof}

This is a leading-order bound, valid once the largest single argument,
$2^{n-1}\mu$, is small and well inside $\arcsin$'s domain -- true from $n=2$
onward at the default $\mu = 2/\|\G\|_S$ (\Cref{eq:pihm-published-prefactor}), but not at $n=1$, where
$\mu \cdot 2^0 = \mu = 1$ exactly and the bound is not yet meaningful.

\subsection{Measured resources and error}
\label{app:published-resources}
% LAND: the paper's claims against the build, n = 2 ... 8: qubits (claimed
%   2n + 3 for G), gates after MCX decomposition (claimed O(n^3)), the measured
%   eps_G(n) beside the bound of prop:epsilon-G rescaled by the published
%   prefactor (eq:pihm-published-prefactor), and alpha / ||G||_2.
% ||G||_S below was recomputed with pihm.operators.derivative on 2026-09-25 and
%   matches the draft's table; the draft's eps_G column used the legacy
%   construction's own rescaling, (N/2)(2/mu), so it is left empty until the
%   Fig. 9 build's values are generated.

\begin{table}[htbp]
    \centering
    \caption{The published encoding of $\G$ against the paper's claims. Classical exact-mode measurement.}
    \label{tab:app-epsilon-G}
    \small
    \begin{tabular}{@{}r r l l l l l@{}}
        \toprule
        $n$ & $N$ & $\|\G\|_S$ & Qubits (claim $2n+3$) & Gates & Bound on $\epsilon_{\G}$ & Measured $\epsilon_{\G}$ \\
        \midrule
        2 & 4   & $\prov{6.15\times10^{1}}$ & \prov{--} & \prov{--} & \prov{--} & \prov{--} \\
        3 & 8   & $\prov{1.18\times10^{3}}$ & \prov{--} & \prov{--} & \prov{--} & \prov{--} \\
        4 & 16  & $\prov{2.04\times10^{4}}$ & \prov{--} & \prov{--} & \prov{--} & \prov{--} \\
        5 & 32  & $\prov{3.38\times10^{5}}$ & \prov{--} & \prov{--} & \prov{--} & \prov{--} \\
        6 & 64  & $\prov{5.50\times10^{6}}$ & \prov{--} & \prov{--} & \prov{--} & \prov{--} \\
        7 & 128 & $\prov{8.88\times10^{7}}$ & \prov{--} & \prov{--} & \prov{--} & \prov{--} \\
        8 & 256 & $\prov{1.43\times10^{9}}$ & \prov{--} & \prov{--} & \prov{--} & \prov{--} \\
        \bottomrule
    \end{tabular}
\end{table}
